% from ega3-1-fr.tex
\oldpage[0\textsubscript{III}]{40}
Supposons que $T$ soit un \emph{foncteur covariant additif} de $\boldsymbol{C}$ dans une
catégorie abélienne $\boldsymbol{C}'$. La définition de
l'\emph{hyperhomologie} $\mathbf{L}_\bullet T(K_\bullet)$ et des \emph{suites
spectrales d'hyperhomologie} de $T$ par rapport à un complexe $K_\bullet$ de
$\boldsymbol{C}$ (lorsqu'elles existent) se fait encore à partir de
(\hyperref[0.11.4.3-fr]{11.4.3}) par «~renversement des flèches~», les termes
$E^2$ des deux suites spectrales ainsi obtenues étant
\[
\label{0.11.6.2.1-fr}
\tag{11.6.2.1}
{}'E^2_{pq}=H_p(L_qT(K_\bullet))
\]
\[
\label{0.11.6.2.2-fr}
\tag{11.6.2.2}
{}''E^2_{pq}=L_pT(H_q(K_\bullet))
\]
où $L_pT$ désigne le $p$\textsuperscript{e} dérivé de $T$ pour $p\geq 0$, et $0$
pour $p<0$, $L_qT(K_\bullet)$ le complexe
$(L_qT(K_i))_{i\in\mathbb Z}$.

Les propriétés de l'hyperhomologie ne se déduisent pas toutes par simple
«~renversement des flèches~» de celles de l'hypercohomologie (à moins qu'on
ne fasse des hypothèses supplémentaires du type AB~5*) de T, 1.5 sur la
catégorie $\boldsymbol{C}'$) en raison des conditions de \emph{régularité} des deux
suites spectrales précédentes, auxquelles il faut cette fois appliquer les
critères de (\hyperref[0.11.3.4-fr]{11.3.4}). Ces derniers montrent que
lorsqu'on suppose que dans $\boldsymbol{C}'$ les limites inductives filtrantes
existent et sont exactes, alors le complexe défini par le bicomplexe
$T(L_{\bullet\bullet})$ existe, et la \emph{seconde} suite spectrale
$\,{}''E(T(L_{\bullet\bullet}))$ est cette fois \emph{régulière}. Si l'on suppose,
soit que $K_\bullet$ soit limité inférieurement, soit qu'il existe un entier
$n$ tel que tout objet de $\boldsymbol{C}$ admette une résolution projective de
longueur $\leq n$, alors les \emph{deux} suites spectrales d'hyperhomologie existent
(sans hypothèse sur $\boldsymbol{C}'$) et sont \emph{birégulières}.
\end{env}

\begin{proposition}[11.6.3]
\label{0.11.6.3-fr}
Soient $\boldsymbol{C}$, $\boldsymbol{C}'$ deux catégories abéliennes, $T$ un foncteur
covariant additif de $\boldsymbol{C}$ dans $\boldsymbol{C}'$. Alors :
\begin{enumerate}
\item[{\rm(i)}] {L'hyperhomologie $\mathbf{L}_\bullet T(K_\bullet)$ est un
foncteur homologique dans la catégorie abélienne des complexes de
$\boldsymbol{C}$ limités inférieurement.}
\item[{\rm(ii)}] {Soit $K_\bullet$ un complexe de $\boldsymbol{C}$ limité
inférieurement. Si $L_nT(K_i)=0$ pour tout $n>0$ et tout $i\in\mathbb Z$, on
a des isomorphismes fonctoriels}
\[
\label{0.11.6.3.1-fr}
\tag{11.6.3.1}
\mathbf{L}_iT(K_\bullet)\xrightarrow{\sim} H_i(T(K_\bullet))\qquad(i\in\mathbb Z).
\]
\item[{\rm(iii)}] {Soit $K_\bullet$ un complexe de $\boldsymbol{C}$ limité
inférieurement. Soit $L_{\bullet\bullet}=(L_{ij})$ un bicomplexe tel que
$L_{ij}=0$ pour $j<0$ et que, pour tout $i$, $L_{i,\bullet}$ soit une
résolution de $K_i$; supposons enfin que $L_nT(L_{ij})=0$ pour tout couple
$(i,j)$ et tout $n>0$. Alors il existe un isomorphisme fonctoriel}
\[
\label{0.11.6.3.2-fr}
\tag{11.6.3.2}
\mathbf{L}_\bullet T(K_\bullet)\xrightarrow{\sim} H_\bullet(T(L_{\bullet\bullet}))
\]
\end{enumerate}
\upshape Les démonstrations se font comme celles de
(\hyperref[0.11.5.2-fr]{11.5.2}), (\hyperref[0.11.4.5-fr]{11.4.5}) et
(\hyperref[0.11.5.3-fr]{11.5.3}) dans le cas des complexes limités
inférieurement. Nous laissons les détails de ces raisonnements au lecteur.
\end{proposition}

\begin{env}[11.6.4]
\label{0.11.6.4-fr}
On a des résultats tout à fait analogues pour les \emph{multifoncteurs} covariants
additifs en chacun des arguments. Par exemple pour un bifoncteur $T$, on a
les deux suites spectrales d'hyperhomologie de termes $E^2$ donnés par
